\documentclass[10pt,a4paper]{amsart}
\usepackage[margin=19mm]{geometry}
\usepackage{amsmath,amssymb,mathtools}
\usepackage[scr=rsfs,cal=euler]{mathalfa}
\usepackage{xfrac}
\usepackage[unicode,hidelinks,bookmarks=false]{hyperref}
\newcommand{\ie}{i.e.\ }
\newcommand{\cf}{cf.\ }
\newcommand{\resp}{respectively\ }
\newcommand{\Subsection}{\subsection}
\providecommand{\nobreakdash}{\nobreak\hbox{-}}
\setlength{\emergencystretch}{2em}
\multlinegap0pt
\usepackage{multirow}
\usepackage{amssymb}
\usepackage{xcolor}
\usepackage{booktabs}
\usepackage{algorithm}\usepackage{algpseudocode}
\newtheorem{definition}{Definition}
\newtheorem{lemma}{Lemma}
\title{Dynamical stability of various convex graphical translators}
\author{Junyoung Park}
\date{Source: arXiv:2601.22368v2 (CC BY 4.0)}
\begin{document}
\maketitle

\noindent\emph{Source and licence.} Dynamical stability of various convex graphical translators. Version arXiv:2601.22368v2 (CC BY 4.0), distributed by arXiv under the Creative Commons Attribution 4.0 International licence, http://creativecommons.org/licenses/by/4.0/, as recorded in the arXiv licence metadata for this exact version and shown on the abstract page https://arxiv.org/abs/2601.22368v2. Full licence notice: \texttt{licenses/arxiv-2601.22368-CC-BY-4.0.txt}.\par\smallskip

\noindent\textit{Editorial scope.} Counted unit: the article's lemma estimate (source0.tex lines 494--507). The article's complete original proof of it is delivered with it (source0.tex lines 508--559, one \texttt{proof} environment of 52 lines). No mathematics is written, repaired or re-derived: the statement and the proof are the article's own lines, in the article's own notation, and no retained line is substituted or re-flowed.  Retained uncounted as the source-local context the proof needs: lines 348--350; lines 382--389; lines 485--490.  Nothing was omitted from any retained span, so the package carries no omission record.  No retained line contains a citation, so the wrapper carries no bibliography and no bibliography entry.  No retained line contains a figure environment, an image-inclusion command or drawing code, so no image is delivered and the package declares no asset dependency.  Editorial formatting only, in addition: an AMS \texttt{amsart} preamble that restates the article's own declarations and macros used by the retained lines, namely the environments \texttt{definition}, \texttt{lemma} on separate counters, together with the standard packages those lines need.  The retained environments print in this document as Definition 1, Lemma 1 in order of appearance, whereas the article prints the same items with the numbering of its own full document.  The complete native source of the article as distributed in the arXiv e-print for this exact version is bundled unchanged as \texttt{sources/193650/source0.tex}.
    \begin{equation}\label{definition of domain}
    K_{r, \lambda} = B_{\mathbf{R}^{n-1}}(0, r) \times (-b + \lambda, b - \lambda) \subset \Omega^n_b.
    \end{equation}

\begin{definition}[Ancient pancake solution]\label{def of ancient pancake}
For each $n \geq 1$, there exists a compact, convex, $O(1)\times O(n)$ symmetric ancient solution to mean curvature flow which lies in the slab
$\Omega = \{x \in \mathbf{R}^{n+1}\ 
 |\ |x_n| < \frac{\pi}{2}\}$, denoted by $\{\Sigma^n_t\}_{t \in (-\infty, 0)}$. Furthermore, as $t \to -\infty$, we have the following `radius' estimate
\begin{align}\label{model pancake asymptotic data}
   & \min_{p \in \Sigma_t}|p| = |P(e_n, t)| \geq \frac{\pi}{2} - o(\frac{1}{(-t)^k}) \text{ for any }k > 0, \\ &\max_{p \in \Sigma_t}|p| = |P(\phi, t)| = -t + (n-1)\ln(-t) + C + o(1) \text{ for any unit vector $\phi \in \{e_n\}^{\perp}$.} 
\end{align}
\end{definition}

    \begin{equation}\label{Grahical mcf IVP}
        \begin{cases}
         \frac{\partial u}{\partial t} = \sqrt{1 + |Du|^2}\textit{\emph{div}}(\frac{Du}{\sqrt{1 + |Du|^2}}) \text{ in }\Omega^n_b \times (0,\infty) \\
         u(x, 0) = u_0 \text{ in }\Omega^n_b
        \end{cases}
    \end{equation}

\begin{lemma}[$C^0$ estimate]\label{C^0 estimate}
    Let $u \in C^{\infty}(\Omega^n_b \times (0, \infty)) \cap C^0(\Omega^n_b \times [0, \infty))$ be the solution to the initial value problem \eqref{Grahical mcf IVP}. Then, for each $T > 0$, $r > 0$, $\lambda \in (0,\frac{1}{2}\min(\frac{\pi}{2},b))$, there exists $R = R(n, \lambda, T)$, and $Q = Q(n,\lambda, T)$ so that we have the $C^0$ estimate
 \begin{equation}\label{C0 estimate equation}
     \|u\|_{C^0(K_{r, 2\lambda}\times [0, T])} \leq \|u_0\|_{C^0(K_{ r + R, \lambda})} + Q(n, \lambda, T),
 \end{equation}
 where $K_{ r, \lambda}$ is given by \eqref{definition of domain}.  Moreover, $R$ is given by
\begin{align}\label{radius equation}
    R(n, \lambda,  T) = \frac{\pi (2T + T_0(n))}{2\lambda} + \frac{2(n-1)\lambda}{\pi}\ln (\frac{\pi^2(2T + T_0(n))}{4\lambda^2})+ \frac{2\lambda C(n)}{\pi} + 1,
\end{align}
and $Q$ is given by
\begin{equation}\label{extra c0 quantity estimate}
    Q(n, \lambda, T) = \frac{\pi T}{2\lambda} + \frac{2(n-1)\lambda}{\pi}\ln 2 + 1
\end{equation}
\end{lemma}

\begin{proof}[Proof of lemma \ref{C^0 estimate}]
Let $\{\Sigma^n_t\}_{t\in(-\infty, 0)}$ be the ancient pancake solution given in definition \ref{def of ancient pancake}. Then in view of \eqref{model pancake asymptotic data}, one can find fixed $T_0 = T_0(n) > 0$ so that whenever $t \leq -T_0$, then
\begin{align}\label{model pancake radius estimate}
     0 < \frac{\pi}{2}- \min_{p \in \Sigma_t}|p| < \frac{1}{4} , \ \ \ |\max_{p \in \Sigma_t}|p| - (-t + (n-1)\ln(-t) + C)| < \frac{1}{4}.
\end{align}
Then, for each $\lambda \in (0,\frac{1}{2}\min(\frac{\pi}{2}, b))$, we can consider the parabolically rescaled ancient pancakes
\begin{equation}\label{definition of rescaled ancient pancakes}
    \Sigma^{n, \lambda}_t = \frac{2\lambda}{\pi} \Sigma^n_{(\frac{2\lambda}{\pi})^{-2}t}.
\end{equation}
Set 
\begin{equation}\label{def of long and short radius}
    r_s(t,\lambda) = \min_{p \in \Sigma^{n, \lambda}_t}|p|, \ \ \ r_l(t, \lambda) =  \max_{p \in \Sigma^{n, \lambda}_t}|p|.
\end{equation}
The inequalities \eqref{model pancake radius estimate} with $\lambda \leq \frac{\pi}{2}$ imply that whenever $t \leq -T_0(n)$, then $(\frac{2\lambda}{\pi})^{-2}t \leq t \leq -T_0(n)$, hence
\begin{equation}\label{short radius estimate for rescaled ancient pancakes}
    0 < \lambda - r_s(t,\lambda)   < \frac{\lambda}{2\pi} \leq \frac{1}{4},
\end{equation}
and
\begin{equation}\label{long radius estimate for rescaled ancient pancakes}
    |r_l(t, \lambda) - (-\frac{\pi t}{2\lambda}+ \frac{2(n-1)\lambda}{\pi}\ln(-\frac{\pi^2t}{4\lambda^2}) + \frac{2\lambda C}{\pi})| \leq \frac{\lambda}{2\pi} \leq \frac{1}{4}.
\end{equation}
We now prove lemma \ref{C^0 estimate}. Let $u$ be given as in lemma \ref{C^0 estimate}. Fix any $x_0 \in K_{ r, 2\lambda}$, and $T > 0$. We consider the $-(T_0 + 2T)$-time slice of translated rescaled ancient pancakes
\begin{equation}\label{def of upper barriers}
    M^{n, \lambda, x_0, +}_{-T_0 - 2T } = \Sigma^{n, \lambda}_{-T_0 -2T}  + x_0 + (\|u_0\|_{C^0(K_{ r + R, \lambda})} + r_l(-T_0-2T, \lambda) + \frac{1}{2})e_{n+1},
\end{equation}
and
\begin{equation}\label{def of lower barriers}
    M^{n, \lambda, x_0, -}_{-T_0 - 2T } = \Sigma^{n, \lambda}_{-T_0 -2T}  + x_0 - (\|u_0\|_{C^0(K_{r + R, \lambda})} + r_l(-T_0-2T, \lambda) + \frac{1}{2})e_{n+1},
\end{equation}
where $r_l$ is given by \eqref{def of long and short radius}. 
Then by definitions \eqref{def of upper barriers}, \eqref{def of lower barriers}, inequalities \eqref{short radius estimate for rescaled ancient pancakes}, \eqref{long radius estimate for rescaled ancient pancakes}, and the definition of $R$ in \eqref{radius equation}, we see that
\begin{equation}\label{good boundary behavior}
    M^{n, \lambda, x_0,+}_{-T_0 - 2T},M^{n, \lambda, x_0,-}_{-T_0 - 2T} \subset K_{r + R, \lambda} \times \mathbf{R},
\end{equation}
and 
\begin{equation}\label{initial avoidance}
     M^{n, \lambda, x_0,+}_{-T_0 - 2T},M^{n, \lambda, x_0,-}_{-T_0 - 2T}, \text{ and graph of }u_0 \text{ are pairwise disjoint}.
\end{equation}
We now evolve $M^{n, \lambda, x_0,+}_{-T_0 - 2T},M^{n, \lambda, x_0,-}_{-T_0 - 2T}, \text{ and graph of }u_0$ by mean curvature flow up to time $t = T$. Due to \eqref{good boundary behavior} and \eqref{initial avoidance}, we can conclude by avoidance principle that  
\begin{equation}\label{later avoidance}
     M^{n, \lambda, x_0,+}_{-T_0 - 2T + t},M^{n, \lambda, x_0,-}_{-T_0 - 2T+ t}, \text{ and graph of }u(\cdot, t) \text{ are pairwise disjoint for all }0 \leq t \leq T.
\end{equation}
Note that by \eqref{def of upper barriers}, \eqref{def of lower barriers} and the symmetry of ancient pancake solution, the centers of $M^{n, \lambda, x_0,+}_{-T_0 - 2T + t},M^{n, \lambda, x_0,-}_{-T_0 - 2T+ t}$ lie on the line $\{x_0\} \times \mathbf{R}$. Then by \eqref{later avoidance}, we can obtain the estimate
\begin{equation}\label{C0 bound}
    |u(x_0, t)| \leq \|u_0\|_{C^0(K_{ r + R, \lambda})} + \frac{1}{2} + r_l(-T_0 - 2T, \lambda) - r_l(-T_0 -T, \lambda)
\end{equation}
Because $-T_0 - 2T, -T_0 - T \leq -T_0$, we can apply \eqref{long radius estimate for rescaled ancient pancakes}. This gives us 
\begin{align*}
        \frac{1}{2} + r_l(-T_0 - 2T, \lambda) &- r_l(-T_0 -T, \lambda)  \\ \leq &\frac{1}{2} + [(\frac{\pi (T_0 + 2T)}{2\lambda}+ \frac{2(n-1)\lambda}{\pi}\ln(\frac{\pi^2(T_0 + 2T)}{4\lambda^2}) + \frac{2\lambda C}{\pi}) + \frac{1}{4}]\\ & - [(\frac{\pi (T_0 + T)}{2\lambda}+ \frac{2(n-1)\lambda}{\pi}\ln(\frac{\pi^2(T_0 + T)}{4\lambda^2}) + \frac{2\lambda C}{\pi}) - \frac{1}{4}]\\ \leq & \frac{\pi T}{2\lambda} + \frac{2(n-1)\lambda}{\pi}\ln 2 + 1
\end{align*}
Combining with \eqref{C0 bound}, and recalling the definition of $Q$ given by \eqref{extra c0 quantity estimate}, we finally have the desired estimate \eqref{C0 estimate equation}.
\end{proof}

\end{document}
